%==============================================================
%  CURSO DE SQL — EXAMEN PRÁCTICO
%  Formato: Opción múltiple, completar código, relacionar
%  Bases: tienda, empleados, ventas, jardinería, universidad
%==============================================================

\documentclass[12pt, a4paper]{article}

%--- Idioma y codificación
\usepackage[utf8]{inputenc}
\usepackage[T1]{fontenc}
\usepackage[spanish, es-noshorthands]{babel}
\usepackage{lmodern}
\usepackage{microtype}
\usepackage{csquotes}

%--- Geometría
\usepackage[top=2.5cm, bottom=2.5cm, left=2.8cm, right=2.8cm]{geometry}

%--- Color
\usepackage[dvipsnames, table]{xcolor}
\definecolor{sqlblue}{RGB}{15,52,96}
\definecolor{sqlaccent}{RGB}{0,150,199}
\definecolor{codebg}{RGB}{245,247,250}
\definecolor{correctbg}{RGB}{232,245,233}
\definecolor{warningbg}{RGB}{255,248,225}
\definecolor{fillcolor}{RGB}{200,220,240}

%--- Cajas
\usepackage{tcolorbox}
\tcbuselibrary{breakable, skins}

\newtcolorbox{ejercicio}[2]{
  enhanced, breakable,
  colback=white, colframe=sqlblue!65,
  coltitle=white, fonttitle=\bfseries\small,
  title={#1 \hfill #2},
  attach boxed title to top left={yshift=-2mm, xshift=4mm},
  boxed title style={colback=sqlblue!80, colframe=sqlblue, sharp corners},
  top=6pt, bottom=6pt, left=6pt, right=6pt
}

\newtcolorbox{solucionbox}[1]{
  enhanced, breakable,
  colback=correctbg, colframe=ForestGreen!55,
  coltitle=white, fonttitle=\bfseries\small,
  title={Solución — #1},
  attach boxed title to top left={yshift=-2mm, xshift=4mm},
  boxed title style={colback=ForestGreen!60, colframe=ForestGreen, sharp corners},
  top=5pt, bottom=5pt, left=6pt, right=6pt
}

\newtcolorbox{hint}{
  colback=warningbg, colframe=Goldenrod!70,
  fonttitle=\bfseries\small, title={Pista},
  attach boxed title to top left={yshift=-2mm, xshift=4mm},
  boxed title style={colback=Goldenrod!65, colframe=Goldenrod, sharp corners},
  top=4pt, bottom=4pt, left=6pt, right=6pt
}

\newtcolorbox{nota}{
  colback=sqlaccent!8, colframe=sqlaccent!40,
  fontupper=\small\itshape,
  left=4pt, right=4pt, top=3pt, bottom=3pt
}

%--- Código SQL
\usepackage{listings}
\lstdefinestyle{sql}{
  language=SQL,
  basicstyle=\ttfamily\small,
  keywordstyle=\bfseries\color{sqlblue},
  stringstyle=\color{BrickRed},
  commentstyle=\itshape\color{gray},
  backgroundcolor=\color{codebg},
  frame=single, framerule=0.4pt, rulecolor=\color{gray!40},
  numbers=left, numberstyle=\tiny\color{gray!60}, numbersep=7pt,
  tabsize=4, showstringspaces=false,
  breaklines=true, xleftmargin=0.5cm,
  literate=
    {á}{{\'{a}}}1 {é}{{\'{e}}}1 {í}{{\'{i}}}1
    {ó}{{\'{o}}}1 {ú}{{\'{u}}}1
    {Á}{{\'{A}}}1 {É}{{\'{E}}}1 {Í}{{\'{I}}}1
    {Ó}{{\'{O}}}1 {Ú}{{\'{U}}}1
    {ñ}{{\~{n}}}1 {Ñ}{{\~{N}}}1,
  morekeywords={
    DATABASE,DATABASES,USE,ENGINE,CHARSET,AUTO_INCREMENT,
    UNSIGNED,VARCHAR,DOUBLE,FLOAT,NUMERIC,TINYINT,SMALLINT,
    DATE,TEXT,ENUM,UNIQUE,FOREIGN,KEY,REFERENCES,ON,DELETE,
    UPDATE,CASCADE,RESTRICT,INNER,LEFT,RIGHT,JOIN,NATURAL,
    BETWEEN,LIKE,IN,NOT,EXISTS,ALL,ANY,HAVING,GROUP,ORDER,
    LIMIT,OFFSET,UNION,DISTINCT,AS,CASE,WHEN,THEN,ELSE,END,
    COUNT,SUM,AVG,MIN,MAX,ROUND,TRUNCATE,UPPER,LOWER,LENGTH,
    SUBSTRING,CONCAT,IFNULL,COALESCE,YEAR,MONTH,DATE_FORMAT,
    ADDDATE,DATEDIFF,NOW,CURDATE,FORMAT,IF,MATCH,AGAINST
  }
}
\lstset{style=sql}

% Blank para completar
\newcommand{\blank}[1][4cm]{\underline{\hspace{#1}}}

%--- Listas
\usepackage{enumitem}
\newlist{opciones}{enumerate}{1}
\setlist[opciones]{label=\textbf{\alph*)}, leftmargin=1.8em, itemsep=4pt, topsep=5pt}

%--- Estrellas de dificultad
\usepackage{fontawesome5}
\newcommand{\estrella}{{\color{Goldenrod}\faStarSolid}}
\newcommand{\estrellavacia}{{\color{gray!45}\faStarSolid}}
\newcommand{\nivel}[1]{%
  \ifnum#1=1\estrella\estrellavacia\estrellavacia\fi
  \ifnum#1=2\estrella\estrella\estrellavacia\fi
  \ifnum#1=3\estrella\estrella\estrella\fi
}

%--- Contador de ejercicios global
\newcounter{numej}
\setcounter{numej}{0}
\newcommand{\ej}{\stepcounter{numej}\thenumej}

%--- Cabecera / pie
\usepackage{fancyhdr}
\pagestyle{fancy}
\fancyhf{}
\fancyhead[L]{\small\textit{Curso de SQL — Examen Práctico}}
\fancyhead[R]{\small\textit{Bases de datos: tienda, empleados, ventas}}
\fancyfoot[C]{\small\thepage}
\renewcommand{\headrulewidth}{0.4pt}
\setlength{\headheight}{14pt}

%--- Tabla de verificación
\usepackage{booktabs, tabularx, multirow, array}

%--- Hipervínculos
\usepackage[colorlinks=true, linkcolor=sqlblue, urlcolor=sqlblue]{hyperref}

%--- Espaciado
\setlength{\parskip}{5pt}
\setlength{\parindent}{0pt}

%==============================================================
\begin{document}
%==============================================================

%--- Portada
\begin{center}
  {\large\textsc{Curso de SQL y Bases de Datos Relacionales}}\\[6pt]
  {\LARGE\textbf{Examen Práctico de Consultas SQL}}\\[4pt]
  {\large Tienda de Informática · Empleados · Ventas}\\[8pt]
  \small
  \begin{tabular}{ll}
    \textbf{Nombre:} & \blank[7cm] \\[6pt]
    \textbf{Fecha:}  & \blank[3cm] \\[6pt]
    \textbf{Grupo:}  & \blank[3cm] \\
  \end{tabular}\\[10pt]
  \hrule
\end{center}

\begin{nota}
\textbf{Instrucciones.}
Cada ejercicio indica su nivel de dificultad:
\nivel{1}~básico · \nivel{2}~intermedio · \nivel{3}~avanzado.
En los ejercicios de opción múltiple solo hay \textbf{una} respuesta correcta.
En los ejercicios de código, escribe en los espacios marcados con \texttt{\_\_\_\_\_\_}.
Las soluciones se encuentran al final del documento.
\end{nota}

\tableofcontents
\newpage

%==============================================================
\section{Base de datos: Tienda de Informática}
%==============================================================

\begin{nota}
Se trabaja con las tablas \texttt{fabricante(id, nombre)} y
\texttt{producto(id, nombre, precio, id\_fabricante)}.
\end{nota}

%--------------------------------------------------------------
\subsection{Consultas sobre una tabla — Opción múltiple}
%--------------------------------------------------------------

\begin{ejercicio}{Ejercicio \ej}{\nivel{1}}
¿Cuál de las siguientes consultas lista \textbf{únicamente} los nombres
de todos los productos de la tabla \texttt{producto}?

\begin{opciones}
  \item \lstinline|SELECT * FROM producto;|
  \item \lstinline|SELECT nombre FROM producto;|
  \item \lstinline|SELECT nombre, precio FROM producto;|
  \item \lstinline|SELECT producto FROM nombre;|
\end{opciones}
\end{ejercicio}

\begin{ejercicio}{Ejercicio \ej}{\nivel{1}}
Se quiere mostrar el nombre de los productos con su precio en
\textbf{céntimos} (multiplicar por 100) usando el alias
\texttt{centimos}. ¿Qué consulta es correcta?

\begin{opciones}
  \item \lstinline|SELECT nombre, precio * 100 FROM producto;|
  \item \lstinline|SELECT nombre, precio / 100 AS centimos FROM producto;|
  \item \lstinline|SELECT nombre, precio * 100 AS centimos FROM producto;|
  \item \lstinline|SELECT nombre AS centimos, precio * 100 FROM producto;|
\end{opciones}
\end{ejercicio}

\begin{ejercicio}{Ejercicio \ej}{\nivel{1}}
¿Cuál de estas consultas devuelve los fabricantes cuyo nombre
\textbf{empieza} por la letra \texttt{S}?

\begin{opciones}
  \item \lstinline|SELECT nombre FROM fabricante WHERE nombre = 'S%';|
  \item \lstinline|SELECT nombre FROM fabricante WHERE nombre LIKE 'S%';|
  \item \lstinline|SELECT nombre FROM fabricante WHERE nombre LIKE '%S';|
  \item \lstinline|SELECT nombre FROM fabricante WHERE nombre IN ('S');|
\end{opciones}
\end{ejercicio}

\begin{ejercicio}{Ejercicio \ej}{\nivel{1}}
¿Cuál es el resultado de la siguiente consulta?
\begin{lstlisting}
SELECT nombre FROM fabricante ORDER BY nombre DESC LIMIT 2 OFFSET 3;
\end{lstlisting}

\begin{opciones}
  \item Los 2 primeros fabricantes en orden descendente.
  \item Los fabricantes en posiciones 4 y 5 del orden descendente.
  \item Los 3 últimos fabricantes en orden ascendente.
  \item Los fabricantes en posiciones 1 y 2 en orden ascendente.
\end{opciones}
\end{ejercicio}

\begin{ejercicio}{Ejercicio \ej}{\nivel{2}}
Se desea listar productos con precio \textbf{entre 60\,€ y 200\,€}
usando \texttt{BETWEEN}. ¿Qué consulta es equivalente a
\lstinline|WHERE precio >= 60 AND precio <= 200|?

\begin{opciones}
  \item \lstinline|WHERE precio BETWEEN 60 OR 200|
  \item \lstinline|WHERE precio BETWEEN 60 AND 200|
  \item \lstinline|WHERE precio IN (60, 200)|
  \item \lstinline|WHERE precio > 60 AND precio < 200|
\end{opciones}
\end{ejercicio}

\begin{ejercicio}{Ejercicio \ej}{\nivel{2}}
Se necesita listar los productos donde el fabricante sea 1, 3 o 5
usando el operador \texttt{IN}. ¿Cuál es la consulta correcta?

\begin{opciones}
  \item \lstinline|WHERE id_fabricante = 1 OR 3 OR 5|
  \item \lstinline|WHERE id_fabricante BETWEEN 1 AND 5|
  \item \lstinline|WHERE id_fabricante IN (1, 3, 5)|
  \item \lstinline|WHERE id_fabricante LIKE (1, 3, 5)|
\end{opciones}
\end{ejercicio}

\begin{ejercicio}{Ejercicio \ej}{\nivel{2}}
¿Cuál de las siguientes consultas devuelve el nombre y precio del
\textbf{producto más barato} usando solo \texttt{ORDER BY} y \texttt{LIMIT}?

\begin{opciones}
  \item \lstinline|SELECT nombre, precio FROM producto ORDER BY precio LIMIT 1;|
  \item \lstinline|SELECT nombre, MIN(precio) FROM producto;|
  \item \lstinline|SELECT nombre, precio FROM producto ORDER BY precio DESC LIMIT 1;|
  \item \lstinline|SELECT nombre, precio FROM producto WHERE precio = MIN(precio);|
\end{opciones}
\end{ejercicio}

%--------------------------------------------------------------
\subsection{Consultas sobre una tabla — Completar código}
%--------------------------------------------------------------

\begin{ejercicio}{Ejercicio \ej}{\nivel{1}}
Completa la consulta para listar los nombres de los fabricantes
en orden \textbf{descendente}:

\begin{lstlisting}
SELECT nombre
FROM   fabricante
ORDER BY ________ ________;
\end{lstlisting}
\end{ejercicio}

\begin{ejercicio}{Ejercicio \ej}{\nivel{1}}
Completa para mostrar nombre y precio de los productos cuyo precio
sea \textbf{mayor o igual a 400\,€}:

\begin{lstlisting}
SELECT nombre, precio
FROM   producto
WHERE  ________ ________ 400;
\end{lstlisting}
\end{ejercicio}

\begin{ejercicio}{Ejercicio \ej}{\nivel{2}}
Completa para obtener los fabricantes cuyo nombre tiene \textbf{exactamente
4 caracteres}:

\begin{lstlisting}
SELECT nombre
FROM   fabricante
WHERE  nombre LIKE '________';
\end{lstlisting}
\end{ejercicio}

\begin{ejercicio}{Ejercicio \ej}{\nivel{2}}
Completa la consulta para mostrar las \textbf{5 primeras filas} de
la tabla \texttt{fabricante}:

\begin{lstlisting}
SELECT *
FROM   fabricante
________ ________;
\end{lstlisting}
\end{ejercicio}

\begin{ejercicio}{Ejercicio \ej}{\nivel{2}}
Completa la consulta para listar los nombres de los productos
convirtiendo el nombre a \textbf{mayúsculas}:

\begin{lstlisting}
SELECT ________(nombre), precio
FROM   producto;
\end{lstlisting}
\end{ejercicio}

\begin{ejercicio}{Ejercicio \ej}{\nivel{2}}
Completa para obtener 2 filas a partir de la \textbf{cuarta fila}
de \texttt{fabricante} (la cuarta fila incluida):

\begin{lstlisting}
SELECT *
FROM   fabricante
LIMIT  ________ OFFSET ________;
\end{lstlisting}
\end{ejercicio}

%--------------------------------------------------------------
\subsection{Consultas multitabla — Opción múltiple}
%--------------------------------------------------------------

\begin{ejercicio}{Ejercicio \ej}{\nivel{2}}
¿Cuál de las siguientes consultas devuelve el nombre del producto,
su precio y el nombre del fabricante para \textbf{todos los productos}?

\begin{opciones}
  \item
\begin{lstlisting}
SELECT p.nombre, p.precio, f.nombre
FROM producto p, fabricante f;
\end{lstlisting}
  \item
\begin{lstlisting}
SELECT p.nombre, p.precio, f.nombre
FROM producto p
INNER JOIN fabricante f ON p.id_fabricante = f.id;
\end{lstlisting}
  \item
\begin{lstlisting}
SELECT p.nombre, p.precio, f.nombre
FROM producto p
LEFT JOIN fabricante f ON p.id = f.id;
\end{lstlisting}
  \item Las opciones a) y b) son equivalentes y correctas.
\end{opciones}
\end{ejercicio}

\begin{ejercicio}{Ejercicio \ej}{\nivel{2}}
Para obtener todos los fabricantes aunque \textbf{no tengan productos}
asociados, ¿qué tipo de JOIN se debe usar?

\begin{opciones}
  \item \texttt{INNER JOIN} con la tabla \texttt{fabricante} a la izquierda.
  \item \texttt{LEFT JOIN} con \texttt{fabricante} como tabla izquierda.
  \item \texttt{RIGHT JOIN} con \texttt{producto} como tabla izquierda.
  \item \texttt{CROSS JOIN} entre ambas tablas.
\end{opciones}
\end{ejercicio}

\begin{ejercicio}{Ejercicio \ej}{\nivel{3}}
¿Cuál de estas consultas devuelve \textbf{solamente los fabricantes
que NO tienen productos} asociados?

\begin{opciones}
  \item
\begin{lstlisting}
SELECT f.nombre FROM fabricante f
INNER JOIN producto p ON f.id = p.id_fabricante
WHERE p.id IS NULL;
\end{lstlisting}
  \item
\begin{lstlisting}
SELECT f.nombre FROM fabricante f
LEFT JOIN producto p ON f.id = p.id_fabricante
WHERE p.id IS NULL;
\end{lstlisting}
  \item
\begin{lstlisting}
SELECT f.nombre FROM fabricante f
WHERE f.id NOT IN (SELECT id FROM fabricante);
\end{lstlisting}
  \item
\begin{lstlisting}
SELECT f.nombre FROM fabricante f
LEFT JOIN producto p ON f.id = p.id_fabricante
WHERE p.id IS NOT NULL;
\end{lstlisting}
\end{opciones}
\end{ejercicio}

%--------------------------------------------------------------
\subsection{Consultas multitabla — Completar código}
%--------------------------------------------------------------

\begin{ejercicio}{Ejercicio \ej}{\nivel{2}}
Completa la consulta para obtener el nombre del producto, precio y
nombre de fabricante de \textbf{todos los productos del fabricante Lenovo}
(usando \texttt{INNER JOIN}):

\begin{lstlisting}
SELECT p.nombre, p.precio, f.nombre
FROM   producto p
________ ________ fabricante f ON p.id_fabricante = f.id
WHERE  f.nombre = ________;
\end{lstlisting}
\end{ejercicio}

\begin{ejercicio}{Ejercicio \ej}{\nivel{2}}
Completa la consulta para listar los fabricantes con sus productos,
incluyendo fabricantes \textbf{sin productos}:

\begin{lstlisting}
SELECT f.nombre AS fabricante,
       p.nombre AS producto
FROM   fabricante f
________ ________ producto p ON f.id = p.id_fabricante;
\end{lstlisting}
\end{ejercicio}

\begin{ejercicio}{Ejercicio \ej}{\nivel{3}}
Completa la consulta para obtener el nombre del producto más caro
del fabricante \textbf{Lenovo} usando una subconsulta:

\begin{lstlisting}
SELECT nombre
FROM   producto
WHERE  precio = (
    SELECT ________(precio)
    FROM   producto
    WHERE  id_fabricante = (
        SELECT ________ FROM fabricante WHERE nombre = 'Lenovo'
    )
);
\end{lstlisting}

\begin{hint}
Piensa qué función de agregación devuelve el valor máximo.
La subconsulta interna recupera el \texttt{id} del fabricante por su nombre.
\end{hint}
\end{ejercicio}

%--------------------------------------------------------------
\subsection{Consultas resumen — Opción múltiple}
%--------------------------------------------------------------

\begin{ejercicio}{Ejercicio \ej}{\nivel{2}}
¿Cuál de las siguientes consultas calcula correctamente el
\textbf{número de fabricantes} distintos que tienen productos en la tabla
\texttt{producto}?

\begin{opciones}
  \item \lstinline|SELECT COUNT(*) FROM producto;|
  \item \lstinline|SELECT COUNT(id_fabricante) FROM producto;|
  \item \lstinline|SELECT COUNT(DISTINCT id_fabricante) FROM producto;|
  \item \lstinline|SELECT COUNT(DISTINCT nombre) FROM producto;|
\end{opciones}
\end{ejercicio}

\begin{ejercicio}{Ejercicio \ej}{\nivel{2}}
Se quiere mostrar el nombre de cada fabricante y el número de productos
que tiene, \textbf{incluyendo fabricantes sin productos} (con valor 0).
¿Cuál es la estructura correcta?

\begin{opciones}
  \item
\begin{lstlisting}
SELECT f.nombre, COUNT(p.id) AS total
FROM   fabricante f
INNER JOIN producto p ON f.id = p.id_fabricante
GROUP BY f.nombre ORDER BY total DESC;
\end{lstlisting}
  \item
\begin{lstlisting}
SELECT f.nombre, COUNT(p.id) AS total
FROM   fabricante f
LEFT JOIN producto p ON f.id = p.id_fabricante
GROUP BY f.nombre ORDER BY total DESC;
\end{lstlisting}
  \item
\begin{lstlisting}
SELECT f.nombre, COUNT(*) AS total
FROM   fabricante f, producto p
GROUP BY f.nombre;
\end{lstlisting}
  \item
\begin{lstlisting}
SELECT f.nombre, SUM(p.id) AS total
FROM   fabricante f
LEFT JOIN producto p ON f.id = p.id_fabricante
GROUP BY f.nombre;
\end{lstlisting}
\end{opciones}
\end{ejercicio}

\begin{ejercicio}{Ejercicio \ej}{\nivel{3}}
¿Cuál de las consultas muestra el \textbf{precio máximo, mínimo y medio}
de los fabricantes cuyo precio medio supera los \textbf{200\,€},
mostrando también el nombre del fabricante?

\begin{opciones}
  \item
\begin{lstlisting}
SELECT f.nombre, MAX(p.precio), MIN(p.precio), AVG(p.precio)
FROM   producto p
INNER JOIN fabricante f ON p.id_fabricante = f.id
WHERE  AVG(p.precio) > 200
GROUP BY f.nombre;
\end{lstlisting}
  \item
\begin{lstlisting}
SELECT f.nombre, MAX(p.precio), MIN(p.precio), AVG(p.precio)
FROM   producto p
INNER JOIN fabricante f ON p.id_fabricante = f.id
GROUP BY f.nombre
HAVING AVG(p.precio) > 200;
\end{lstlisting}
  \item
\begin{lstlisting}
SELECT f.nombre, MAX(p.precio), MIN(p.precio), AVG(p.precio)
FROM   producto p
INNER JOIN fabricante f ON p.id_fabricante = f.id
GROUP BY f.nombre
HAVING COUNT(*) > 200;
\end{lstlisting}
  \item Las opciones a) y b) son equivalentes.
\end{opciones}
\end{ejercicio}

%--------------------------------------------------------------
\subsection{Subconsultas — Completar código}
%--------------------------------------------------------------

\begin{ejercicio}{Ejercicio \ej}{\nivel{2}}
Completa para obtener todos los productos del fabricante \textbf{Lenovo}
\emph{sin usar} \texttt{INNER JOIN}:

\begin{lstlisting}
SELECT *
FROM   producto
WHERE  id_fabricante = (
    SELECT ________
    FROM   fabricante
    WHERE  nombre = 'Lenovo'
);
\end{lstlisting}
\end{ejercicio}

\begin{ejercicio}{Ejercicio \ej}{\nivel{3}}
Completa la subconsulta para obtener el nombre del producto más caro
de toda la tabla, \textbf{sin usar} \texttt{MAX}, \texttt{ORDER BY}
ni \texttt{LIMIT} (usando \texttt{ALL}):

\begin{lstlisting}
SELECT nombre, precio
FROM   producto
WHERE  precio >= ________ (
    SELECT precio FROM producto
);
\end{lstlisting}

\begin{hint}
\texttt{ALL} devuelve verdadero si la condición se cumple para
\emph{todos} los valores de la subconsulta.
Si el precio es $\geq$ que \emph{todos} los precios, es el máximo.
\end{hint}
\end{ejercicio}

\begin{ejercicio}{Ejercicio \ej}{\nivel{3}}
Completa la consulta para obtener los fabricantes que \textbf{tienen
productos} usando \texttt{EXISTS}:

\begin{lstlisting}
SELECT f.nombre
FROM   fabricante f
WHERE  ________ (
    SELECT *
    FROM   producto p
    WHERE  p.id_fabricante = ________
);
\end{lstlisting}
\end{ejercicio}

%==============================================================
\section{Base de datos: Gestión de Empleados}
%==============================================================

\begin{nota}
Tablas: \texttt{departamento(id, nombre, presupuesto, gastos)} y
\texttt{empleado(id, nif, nombre, apellido1, apellido2, id\_departamento)}.
\end{nota}

%--------------------------------------------------------------
\subsection{Consultas sobre una tabla — Opción múltiple}
%--------------------------------------------------------------

\begin{ejercicio}{Ejercicio \ej}{\nivel{1}}
¿Cuál de las siguientes consultas muestra el nombre y apellidos de
todos los empleados en \textbf{una única columna}?

\begin{opciones}
  \item \lstinline|SELECT nombre, apellido1 FROM empleado;|
  \item \lstinline|SELECT CONCAT(nombre, ' ', apellido1, ' ', apellido2) FROM empleado;|
  \item \lstinline|SELECT nombre + ' ' + apellido1 FROM empleado;|
  \item \lstinline|SELECT UPPER(nombre, apellido1) FROM empleado;|
\end{opciones}
\end{ejercicio}

\begin{ejercicio}{Ejercicio \ej}{\nivel{1}}
Se quiere calcular el \textbf{presupuesto actual} de cada departamento
(presupuesto inicial menos gastos). ¿Cuál es la consulta correcta?

\begin{opciones}
  \item \lstinline|SELECT nombre, presupuesto + gastos AS presupuesto_actual FROM departamento;|
  \item \lstinline|SELECT nombre, presupuesto - gastos AS presupuesto_actual FROM departamento;|
  \item \lstinline|SELECT nombre, gastos - presupuesto AS presupuesto_actual FROM departamento;|
  \item \lstinline|SELECT nombre, presupuesto / gastos AS presupuesto_actual FROM departamento;|
\end{opciones}
\end{ejercicio}

\begin{ejercicio}{Ejercicio \ej}{\nivel{2}}
¿Cuál de las consultas devuelve los empleados cuyo \textbf{segundo
apellido es NULL}?

\begin{opciones}
  \item \lstinline|WHERE apellido2 = NULL|
  \item \lstinline|WHERE apellido2 = ''|
  \item \lstinline|WHERE apellido2 IS NULL|
  \item \lstinline|WHERE ISNULL(apellido2) = 0|
\end{opciones}
\end{ejercicio}

\begin{ejercicio}{Ejercicio \ej}{\nivel{2}}
Se quiere listar los departamentos con \textbf{gastos mayores que el
presupuesto}. ¿Cuál es la consulta adecuada?

\begin{opciones}
  \item \lstinline|WHERE presupuesto > gastos|
  \item \lstinline|WHERE gastos > presupuesto|
  \item \lstinline|WHERE presupuesto = gastos|
  \item \lstinline|WHERE gastos IS NOT NULL|
\end{opciones}
\end{ejercicio}

%--------------------------------------------------------------
\subsection{Consultas sobre una tabla — Completar código}
%--------------------------------------------------------------

\begin{ejercicio}{Ejercicio \ej}{\nivel{1}}
Completa para listar los apellidos de empleados
\textbf{sin repetidos}:

\begin{lstlisting}
SELECT ________ apellido1
FROM   empleado;
\end{lstlisting}
\end{ejercicio}

\begin{ejercicio}{Ejercicio \ej}{\nivel{2}}
Completa para listar los 3 departamentos con \textbf{mayor presupuesto}:

\begin{lstlisting}
SELECT nombre, presupuesto
FROM   departamento
ORDER BY ________ ________
________ 3;
\end{lstlisting}
\end{ejercicio}

\begin{ejercicio}{Ejercicio \ej}{\nivel{2}}
Completa para obtener 5 filas a partir de la \textbf{tercera fila} de
\texttt{empleado} (incluida):

\begin{lstlisting}
SELECT *
FROM   empleado
LIMIT  ________ OFFSET ________;
\end{lstlisting}
\end{ejercicio}

%--------------------------------------------------------------
\subsection{Consultas multitabla — Opción múltiple}
%--------------------------------------------------------------

\begin{ejercicio}{Ejercicio \ej}{\nivel{2}}
¿Cuál de las consultas devuelve un listado con \textbf{todos los empleados},
incluyendo los que \textbf{no tienen departamento} asignado?

\begin{opciones}
  \item
\begin{lstlisting}
SELECT e.*, d.nombre FROM empleado e
INNER JOIN departamento d ON e.id_departamento = d.id;
\end{lstlisting}
  \item
\begin{lstlisting}
SELECT e.*, d.nombre FROM empleado e
LEFT JOIN departamento d ON e.id_departamento = d.id;
\end{lstlisting}
  \item
\begin{lstlisting}
SELECT e.*, d.nombre FROM departamento d
LEFT JOIN empleado e ON e.id_departamento = d.id;
\end{lstlisting}
  \item
\begin{lstlisting}
SELECT e.*, d.nombre FROM empleado e, departamento d
WHERE e.id_departamento = d.id;
\end{lstlisting}
\end{opciones}
\end{ejercicio}

\begin{ejercicio}{Ejercicio \ej}{\nivel{3}}
Se necesita listar los departamentos \textbf{sin empleados} y los
empleados \textbf{sin departamento} en un solo resultado. ¿Cómo se
logra esto en SQL?

\begin{opciones}
  \item Con un único \texttt{LEFT JOIN} de empleado hacia departamento.
  \item Con \texttt{INNER JOIN} y filtrando \texttt{NULL}s.
  \item Con dos consultas unidas por \texttt{UNION}: una con
        \texttt{LEFT JOIN} (empleados sin departamento) y otra con
        \texttt{LEFT JOIN} al revés (departamentos sin empleados).
  \item No es posible en SQL estándar.
\end{opciones}

\begin{hint}
\texttt{UNION} combina resultados de dos \texttt{SELECT} distintos.
Puedes hacer un \texttt{LEFT JOIN} en cada dirección y unirlos.
\end{hint}
\end{ejercicio}

%--------------------------------------------------------------
\subsection{Consultas resumen — Completar código}
%--------------------------------------------------------------

\begin{ejercicio}{Ejercicio \ej}{\nivel{2}}
Completa para calcular el número de empleados en cada departamento,
mostrando el \textbf{nombre} del departamento y el conteo:

\begin{lstlisting}
SELECT d.nombre, COUNT(________)  AS num_empleados
FROM   departamento d
________ ________ empleado e ON e.id_departamento = d.id
GROUP BY ________
ORDER BY num_empleados DESC;
\end{lstlisting}
\end{ejercicio}

\begin{ejercicio}{Ejercicio \ej}{\nivel{3}}
Completa la consulta para mostrar los departamentos con
\textbf{más de 2 empleados}:

\begin{lstlisting}
SELECT d.nombre, COUNT(e.id) AS total
FROM   departamento d
INNER JOIN empleado e ON e.id_departamento = d.id
GROUP BY d.nombre
________ ________(e.id) > 2;
\end{lstlisting}
\end{ejercicio}

%--------------------------------------------------------------
\subsection{Subconsultas — Opción múltiple}
%--------------------------------------------------------------

\begin{ejercicio}{Ejercicio \ej}{\nivel{3}}
¿Cuál de las consultas devuelve el nombre del departamento con
\textbf{mayor presupuesto} sin usar \texttt{MAX}, \texttt{ORDER BY}
ni \texttt{LIMIT} (usando \texttt{ALL})?

\begin{opciones}
  \item
\begin{lstlisting}
SELECT nombre FROM departamento
WHERE presupuesto > ALL (SELECT presupuesto FROM departamento);
\end{lstlisting}
  \item
\begin{lstlisting}
SELECT nombre FROM departamento
WHERE presupuesto >= ALL (SELECT presupuesto FROM departamento);
\end{lstlisting}
  \item
\begin{lstlisting}
SELECT nombre FROM departamento
WHERE presupuesto = ANY (SELECT presupuesto FROM departamento);
\end{lstlisting}
  \item
\begin{lstlisting}
SELECT nombre FROM departamento
WHERE presupuesto IN (SELECT MAX(presupuesto) FROM departamento);
\end{lstlisting}
\end{opciones}
\end{ejercicio}

%==============================================================
\section{Base de datos: Gestión de Ventas}
%==============================================================

\begin{nota}
Tablas: \texttt{cliente}, \texttt{comercial} y
\texttt{pedido(id, total, fecha, id\_cliente, id\_comercial)}.
\end{nota}

%--------------------------------------------------------------
\subsection{Consultas sobre una tabla — Opción múltiple}
%--------------------------------------------------------------

\begin{ejercicio}{Ejercicio \ej}{\nivel{1}}
¿Cuál de las consultas devuelve todos los pedidos realizados durante
el \textbf{año 2017} con total superior a \textbf{500\,€}?

\begin{opciones}
  \item
\begin{lstlisting}
SELECT * FROM pedido WHERE YEAR(fecha) = 2017 AND total > 500;
\end{lstlisting}
  \item
\begin{lstlisting}
SELECT * FROM pedido WHERE fecha = 2017 AND total > 500;
\end{lstlisting}
  \item
\begin{lstlisting}
SELECT * FROM pedido WHERE fecha LIKE '2017%' OR total > 500;
\end{lstlisting}
  \item
\begin{lstlisting}
SELECT * FROM pedido WHERE YEAR(fecha) > 2017 AND total > 500;
\end{lstlisting}
\end{opciones}
\end{ejercicio}

\begin{ejercicio}{Ejercicio \ej}{\nivel{2}}
¿Cuál de las siguientes consultas devuelve los nombres de los
comerciales que \textbf{no} aparecen en la tabla \texttt{pedido}
(usando \texttt{NOT IN})?

\begin{opciones}
  \item
\begin{lstlisting}
SELECT nombre FROM comercial
WHERE id NOT IN (SELECT id FROM pedido);
\end{lstlisting}
  \item
\begin{lstlisting}
SELECT nombre FROM comercial
WHERE id NOT IN (SELECT id_comercial FROM pedido);
\end{lstlisting}
  \item
\begin{lstlisting}
SELECT nombre FROM comercial
WHERE id IN (SELECT id_comercial FROM pedido);
\end{lstlisting}
  \item
\begin{lstlisting}
SELECT nombre FROM comercial
WHERE id IS NULL;
\end{lstlisting}
\end{opciones}
\end{ejercicio}

%--------------------------------------------------------------
\subsection{Consultas sobre una tabla — Completar código}
%--------------------------------------------------------------

\begin{ejercicio}{Ejercicio \ej}{\nivel{1}}
Completa para devolver todos los pedidos ordenados por fecha
(\textbf{más recientes primero}):

\begin{lstlisting}
SELECT *
FROM   pedido
ORDER BY ________ ________;
\end{lstlisting}
\end{ejercicio}

\begin{ejercicio}{Ejercicio \ej}{\nivel{2}}
Completa para listar los clientes que \textbf{no empiezan por A}:

\begin{lstlisting}
SELECT nombre
FROM   cliente
WHERE  nombre ________ LIKE 'A%'
ORDER BY nombre;
\end{lstlisting}
\end{ejercicio}

%--------------------------------------------------------------
\subsection{Consultas multitabla — Completar código}
%--------------------------------------------------------------

\begin{ejercicio}{Ejercicio \ej}{\nivel{2}}
Completa la consulta para listar todos los clientes junto con sus
pedidos, incluyendo clientes \textbf{sin pedidos}:

\begin{lstlisting}
SELECT c.nombre, c.apellido1, p.total, p.fecha
FROM   cliente c
________ ________ pedido p ON c.id = p.id_cliente
ORDER BY c.apellido1, c.apellido2, c.nombre;
\end{lstlisting}
\end{ejercicio}

\begin{ejercicio}{Ejercicio \ej}{\nivel{3}}
Completa la consulta para mostrar todos los pedidos de la clienta
\textbf{Adela Salas Díaz} sin usar \texttt{INNER JOIN}:

\begin{lstlisting}
SELECT *
FROM   pedido
WHERE  id_cliente = (
    SELECT ________
    FROM   cliente
    WHERE  nombre    = 'Adela'
      AND  apellido1 = 'Salas'
      AND  apellido2 = ________
);
\end{lstlisting}
\end{ejercicio}

%--------------------------------------------------------------
\subsection{Consultas resumen — Opción múltiple}
%--------------------------------------------------------------

\begin{ejercicio}{Ejercicio \ej}{\nivel{2}}
¿Cuál de las consultas devuelve el \textbf{número total de pedidos
realizados cada año}?

\begin{opciones}
  \item
\begin{lstlisting}
SELECT fecha, COUNT(*) FROM pedido GROUP BY fecha;
\end{lstlisting}
  \item
\begin{lstlisting}
SELECT YEAR(fecha) AS anio, COUNT(*) AS total
FROM pedido GROUP BY YEAR(fecha) ORDER BY anio;
\end{lstlisting}
  \item
\begin{lstlisting}
SELECT YEAR(fecha) AS anio, SUM(total)
FROM pedido GROUP BY anio;
\end{lstlisting}
  \item
\begin{lstlisting}
SELECT COUNT(*) FROM pedido WHERE YEAR(fecha) > 0;
\end{lstlisting}
\end{opciones}
\end{ejercicio}

\begin{ejercicio}{Ejercicio \ej}{\nivel{3}}
¿Cuál de las consultas devuelve el número total de pedidos
de \textbf{cada cliente}, incluyendo los que \textbf{no han realizado
ningún pedido} (con valor 0)?

\begin{opciones}
  \item
\begin{lstlisting}
SELECT c.id, c.nombre, COUNT(p.id) AS total
FROM cliente c, pedido p
WHERE c.id = p.id_cliente
GROUP BY c.id;
\end{lstlisting}
  \item
\begin{lstlisting}
SELECT c.id, c.nombre, COUNT(p.id) AS total
FROM cliente c
LEFT JOIN pedido p ON c.id = p.id_cliente
GROUP BY c.id, c.nombre;
\end{lstlisting}
  \item
\begin{lstlisting}
SELECT c.id, c.nombre, COUNT(*) AS total
FROM cliente c
LEFT JOIN pedido p ON c.id = p.id_cliente
GROUP BY c.id;
\end{lstlisting}
  \item
\begin{lstlisting}
SELECT c.id, c.nombre, SUM(p.id) AS total
FROM cliente c
LEFT JOIN pedido p ON c.id = p.id_cliente
GROUP BY c.id;
\end{lstlisting}
\end{opciones}

\begin{hint}
Cuando hay \texttt{LEFT JOIN} y se quiere contar solo pedidos
reales (no NULL), usa \texttt{COUNT(p.id)} en lugar de
\texttt{COUNT(*)}.
\end{hint}
\end{ejercicio}

%--------------------------------------------------------------
\subsection{Ejercicio integrador de código completo}
%--------------------------------------------------------------

\begin{ejercicio}{Ejercicio \ej}{\nivel{3}}
Analiza la siguiente consulta e identifica \textbf{cuatro errores}:

\begin{lstlisting}
SELECT c.nombre, COUNT(*) AS num_pedidos,
       AVG(p.total) as promedio
FROM   cliente c,
       pedido p
WHERE  AVG(p.total) > 500
GROUP  BY c.nombre, c.id
ORDER  BY COUNT(*) ASC;
\end{lstlisting}

\vspace{4pt}
\textbf{Error 1:} \blank[9cm]\\[6pt]
\textbf{Error 2:} \blank[9cm]\\[6pt]
\textbf{Error 3:} \blank[9cm]\\[6pt]
\textbf{Error 4:} \blank[9cm]

\begin{hint}
Piensa en: (1) cómo se hace un JOIN entre tablas,
(2) dónde se usan las funciones de agregación para filtrar,
(3) qué columna debería estar en \texttt{GROUP BY},
(4) ¿el orden debería ser ASC o DESC para ver los clientes
con más pedidos primero?
\end{hint}
\end{ejercicio}

%==============================================================
\section{Preguntas de Relacionar Columnas}
%==============================================================

\begin{ejercicio}{Ejercicio \ej}{\nivel{1}}
Relaciona cada cláusula SQL con su descripción correcta:

\vspace{6pt}
\begin{center}
\small
\renewcommand{\arraystretch}{1.5}
\begin{tabular}{@{} p{3.5cm} p{0.8cm} p{7cm} @{}}
  \toprule
  \textbf{Cláusula} & & \textbf{Descripción} \\
  \midrule
  \texttt{WHERE}    & \blank[0.6cm] & A. Filtra grupos después de agrupar \\
  \texttt{HAVING}   & \blank[0.6cm] & B. Elimina filas duplicadas del resultado \\
  \texttt{GROUP BY} & \blank[0.6cm] & C. Filtra filas antes de agrupar \\
  \texttt{DISTINCT} & \blank[0.6cm] & D. Agrupa filas con valores iguales \\
  \texttt{ORDER BY} & \blank[0.6cm] & E. Ordena el resultado \\
  \bottomrule
\end{tabular}
\end{center}
\end{ejercicio}

\begin{ejercicio}{Ejercicio \ej}{\nivel{2}}
Relaciona cada tipo de JOIN con el resultado que produce:

\vspace{6pt}
\begin{center}
\small
\renewcommand{\arraystretch}{1.5}
\begin{tabular}{@{} p{3.5cm} p{0.8cm} p{9cm} @{}}
  \toprule
  \textbf{Tipo JOIN} & & \textbf{Resultado} \\
  \midrule
  \texttt{INNER JOIN}  & \blank[0.6cm] & A. Todas las filas de la tabla derecha + coincidencias \\
  \texttt{LEFT JOIN}   & \blank[0.6cm] & B. Solo filas con coincidencia en ambas tablas \\
  \texttt{RIGHT JOIN}  & \blank[0.6cm] & C. Todas las filas de la tabla izquierda + coincidencias \\
  \texttt{UNION}       & \blank[0.6cm] & D. Resultados de dos SELECT combinados sin duplicados \\
  \texttt{UNION ALL}   & \blank[0.6cm] & E. Resultados de dos SELECT combinados con duplicados \\
  \bottomrule
\end{tabular}
\end{center}
\end{ejercicio}

\begin{ejercicio}{Ejercicio \ej}{\nivel{2}}
Relaciona cada función con su comportamiento frente a valores \texttt{NULL}:

\vspace{6pt}
\begin{center}
\small
\renewcommand{\arraystretch}{1.5}
\begin{tabular}{@{} p{4cm} p{0.8cm} p{8cm} @{}}
  \toprule
  \textbf{Función} & & \textbf{Comportamiento} \\
  \midrule
  \texttt{COUNT(*)}       & \blank[0.6cm] & A. Ignora NULLs; retorna el primer no-NULL de una lista \\
  \texttt{COUNT(columna)} & \blank[0.6cm] & B. Cuenta todas las filas, incluyendo NULLs \\
  \texttt{AVG(columna)}   & \blank[0.6cm] & C. Cuenta solo filas donde la columna no es NULL \\
  \texttt{COALESCE()}     & \blank[0.6cm] & D. Ignora NULLs en el cálculo del promedio \\
  \texttt{IFNULL(a,b)}    & \blank[0.6cm] & E. Devuelve b si a es NULL; equivalente a COALESCE de 2 args \\
  \bottomrule
\end{tabular}
\end{center}
\end{ejercicio}

%==============================================================
\section{Verdadero / Falso}
%==============================================================

\begin{ejercicio}{Ejercicio \ej}{\nivel{1}}
Indica si cada afirmación es \textbf{Verdadera (V)} o
\textbf{Falsa (F)}:

\vspace{4pt}
\begin{enumerate}[label=\arabic*., leftmargin=1.5em, itemsep=8pt]
  \item \blank[1.2cm] \texttt{WHERE} se evalúa después de \texttt{GROUP BY}.
  \item \blank[1.2cm] \texttt{NULL = NULL} devuelve \texttt{TRUE} en SQL.
  \item \blank[1.2cm] \texttt{DISTINCT} elimina filas duplicadas del resultado.
  \item \blank[1.2cm] \texttt{TRUNCATE} admite cláusula \texttt{WHERE}.
  \item \blank[1.2cm] Un \texttt{LEFT JOIN} incluye siempre todas las filas
        de la tabla izquierda aunque no haya coincidencia.
  \item \blank[1.2cm] \texttt{HAVING} puede usarse sin \texttt{GROUP BY}.
  \item \blank[1.2cm] \texttt{COUNT(*)} ignora las filas con valores \texttt{NULL}.
  \item \blank[1.2cm] \texttt{BETWEEN 60 AND 200} incluye los valores 60 y 200.
  \item \blank[1.2cm] \texttt{LIMIT 3 OFFSET 3} devuelve filas 4, 5 y 6.
  \item \blank[1.2cm] \texttt{AVG(columna)} divide entre el total de filas,
        incluyendo las que tienen \texttt{NULL} en esa columna.
\end{enumerate}
\end{ejercicio}

%==============================================================
\newpage
\section*{Soluciones}
\addcontentsline{toc}{section}{Soluciones}
%==============================================================

\begin{solucionbox}{Sección 1 — Tienda de Informática}
\textbf{Opción múltiple:}
Ej.~1~b) \quad Ej.~2~c) \quad Ej.~3~b) \quad Ej.~4~b) \quad
Ej.~5~b) \quad Ej.~6~c) \quad Ej.~7~a)

\medskip
\textbf{Completar código:}
\begin{itemize}[leftmargin=1em, itemsep=3pt]
  \item Ej.~8: \lstinline|ORDER BY nombre DESC|
  \item Ej.~9: \lstinline|WHERE precio >= 400|
  \item Ej.~10: \lstinline|LIKE '____'| (cuatro guiones bajos)
  \item Ej.~11: \lstinline|LIMIT 5|
  \item Ej.~12: \lstinline|UPPER(nombre)|
  \item Ej.~13: \lstinline|LIMIT 2 OFFSET 3|
\end{itemize}

\medskip
\textbf{Multitabla opción múltiple:}
Ej.~14~d) \quad Ej.~15~b) \quad Ej.~16~b)

\medskip
\textbf{Multitabla completar:}
\begin{itemize}[leftmargin=1em, itemsep=3pt]
  \item Ej.~17: \lstinline|INNER JOIN ... 'Lenovo'|
  \item Ej.~18: \lstinline|LEFT JOIN|
  \item Ej.~19: \lstinline|MAX(precio)| ... \lstinline|id|
\end{itemize}

\medskip
\textbf{Resumen opción múltiple:}
Ej.~20~c) \quad Ej.~21~b) \quad Ej.~22~b)

\medskip
\textbf{Subconsultas completar:}
\begin{itemize}[leftmargin=1em, itemsep=3pt]
  \item Ej.~23: \lstinline|id|
  \item Ej.~24: \lstinline|ALL|
  \item Ej.~25: \lstinline|EXISTS| \ldots \lstinline|f.id|
\end{itemize}
\end{solucionbox}

\begin{solucionbox}{Sección 2 — Empleados}
\textbf{Opción múltiple:}
Ej.~26~b) \quad Ej.~27~b) \quad Ej.~28~c) \quad Ej.~29~b)

\medskip
\textbf{Completar código:}
\begin{itemize}[leftmargin=1em, itemsep=3pt]
  \item Ej.~30: \lstinline|DISTINCT|
  \item Ej.~31: \lstinline|presupuesto DESC| \ldots \lstinline|LIMIT|
  \item Ej.~32: \lstinline|LIMIT 5 OFFSET 2|
\end{itemize}

\medskip
\textbf{Multitabla opción múltiple:}
Ej.~33~b) \quad Ej.~34~c)

\medskip
\textbf{Resumen completar:}
\begin{itemize}[leftmargin=1em, itemsep=3pt]
  \item Ej.~35: \lstinline|COUNT(e.id)| \ldots \lstinline|LEFT JOIN| \ldots \lstinline|d.nombre|
  \item Ej.~36: \lstinline|HAVING COUNT(e.id) > 2|
\end{itemize}

\medskip
\textbf{Subconsultas opción múltiple:}
Ej.~37~b)
\end{solucionbox}

\begin{solucionbox}{Sección 3 — Ventas}
\textbf{Opción múltiple:}
Ej.~38~a) \quad Ej.~39~b)

\medskip
\textbf{Completar código:}
\begin{itemize}[leftmargin=1em, itemsep=3pt]
  \item Ej.~40: \lstinline|ORDER BY fecha DESC|
  \item Ej.~41: \lstinline|NOT LIKE 'A%'| (o \lstinline|WHERE nombre NOT LIKE 'A%'|)
  \item Ej.~42: \lstinline|LEFT JOIN|
  \item Ej.~43: \lstinline|id| \ldots \lstinline|'Díaz'|
\end{itemize}

\medskip
\textbf{Resumen opción múltiple:}
Ej.~44~b) \quad Ej.~45~b)

\medskip
\textbf{Ejercicio integrador (Ej.~46) — 4 errores:}
\begin{enumerate}[leftmargin=1.5em]
  \item Falta la condición de JOIN en \texttt{WHERE}:
        \lstinline|WHERE c.id = p.id_cliente|.
  \item \texttt{AVG(p.total) > 500} en \texttt{WHERE} es incorrecto;
        las funciones de agregación van en \texttt{HAVING}.
  \item En \texttt{GROUP BY} solo debe aparecer \texttt{c.id}
        (o \texttt{c.nombre}); no ambos si \texttt{nombre} no es
        clave (o bien agregar \texttt{c.nombre} de forma consistente).
  \item \texttt{ORDER BY COUNT(*) ASC} muestra los menos primero;
        para ver los más frecuentes primero se necesita \texttt{DESC}.
\end{enumerate}
\end{solucionbox}

\begin{solucionbox}{Sección 4 — Relacionar columnas}
\textbf{Ej.~47 — Cláusulas:}
WHERE → C \quad HAVING → A \quad GROUP BY → D \quad
DISTINCT → B \quad ORDER BY → E

\medskip
\textbf{Ej.~48 — JOINs:}
INNER JOIN → B \quad LEFT JOIN → C \quad RIGHT JOIN → A \quad
UNION → D \quad UNION ALL → E

\medskip
\textbf{Ej.~49 — NULLs:}
COUNT(*) → B \quad COUNT(col) → C \quad AVG → D \quad
COALESCE → A \quad IFNULL → E
\end{solucionbox}

\begin{solucionbox}{Sección 5 — Verdadero / Falso (Ej.~50)}
\begin{enumerate}[label=\arabic*., leftmargin=1.5em, itemsep=3pt]
  \item \textbf{F} — \texttt{WHERE} se evalúa \emph{antes} de \texttt{GROUP BY}.
  \item \textbf{F} — \texttt{NULL = NULL} devuelve \texttt{NULL}, no \texttt{TRUE}. Usar \texttt{IS NULL}.
  \item \textbf{V}
  \item \textbf{F} — \texttt{TRUNCATE} no admite \texttt{WHERE}; elimina todas las filas.
  \item \textbf{V}
  \item \textbf{V} — aunque es inusual, es válido sintácticamente.
  \item \textbf{F} — \texttt{COUNT(*)} cuenta \emph{todas} las filas, incluyendo NULLs.
  \item \textbf{V} — \texttt{BETWEEN} es un rango \emph{inclusivo}.
  \item \textbf{V} — \texttt{OFFSET 3} salta 3 filas; devuelve la 4.ª, 5.ª y 6.ª.
  \item \textbf{F} — \texttt{AVG} ignora los NULLs; divide entre el número de valores no nulos.
\end{enumerate}
\end{solucionbox}

\end{document}
